\begin{table*}[tbp]
\centering
\small
\begin{threeparttable}
\caption{Comprehensive candidate policy families. The number of executable candidates is state/variant dependent (8--12 in the separate rolling records).}
\label{tab:candidate_policies}
\begin{tabularx}{0.96\textwidth}{l l X}
\toprule
\rowcolor{TableHeader}
Policy & Main criterion & Engineering interpretation \\
\midrule
\rowcolor{RowWeak}
NP & No protection & Unprotected reference baseline. \\
\rowcolor{RowNeutral}
RP & Random selection & Non-informative defense without cyber or physical knowledge. \\
\rowcolor{RowBaseline}
CDP & Communication degree & Protects communication hubs to suppress malware diffusion. \\
\rowcolor{RowPhysical}
PDP & Power-layer importance & Protects nodes associated with physically important buses. \\
\rowcolor{RowPhysical}
CP & Capacity-proxy protection & Protects nodes coupled to high-capacity or high-impact buses. \\
\rowcolor{RowPhysical}
PHYSICS & Physical-impact score & Protects nodes whose coupled buses have high physics-informed impact scores. \\
\rowcolor{RowScenario}
HYBRID & Cyber--physical weighted score & Balances communication centrality and physical consequence. \\
\rowcolor{RowRisk}
DRAD-score & Dynamic risk-aware score & Uses topology, physical consequence, infection pressure, and information exposure. \\
\rowcolor{RowRisk}
Rank-ensemble & Rank aggregation & Combines cyber, physical, and infection-based rankings. \\
\rowcolor{RowProposed}
CDP-PHYSICS / CDP-PDP & CDP plus physical/PDP supplement & Keeps communication-hub coverage while adding physical-consequence awareness. \\
\bottomrule
\end{tabularx}
\end{threeparttable}
\end{table*}
